\documentclass[11pt]{article}
\usepackage[margin=25mm]{geometry}
\usepackage{amsmath,booktabs,graphicx,hyperref}
\newcommand{\GateStatus}{validation\_pending}
\newcommand{\DevelopmentStars}{9}
\newcommand{\TessSectors}{11}

\title{Conditional utility and transfer limits of topological asteroseismic descriptors}
\author{Gabriel Wendell}
\date{Stage 12 working draft}
\begin{document}\maketitle
\begin{abstract}
We implement a grouped comparison of compact classical and hybrid seismic models,
physical Kepler-to-TESS degradation and frozen-method quality screening. The current
dataset supports \DevelopmentStars{} usable development stars; \TessSectors{} actual
Rudy TESS sectors supply a descriptive transfer case. Scientific G12 is
\GateStatus. No positive population increment, calibrated detection or population
transfer claim follows from this implementation snapshot.
\end{abstract}
\section{Question and controls}
The relevant practical question is whether a small collection of geometric
descriptors improves prediction after conventional seismic estimates and matched
controls are included. Path-superlevel $H_0$ persistence and peak prominence are
the same scalar information. Their location descriptor belongs in both models;
the baseline additionally includes a conventional graph-ridge score. Solar-like
oscillation scaling and observational context follow standard reviews
\cite{ChaplinMiglio2013,GarciaBallot2019}. The hybrid
adds a nominal comb score, the excess competing-alias score and a cylindrical
geometry descriptor. A gain within a compact regression class would establish
conditional operational utility, not independent spectral information.
The observations and published reference centers follow the LEGACY studies
\cite{Lund2017,SilvaAguirre2017,Lund2017Erratum}.
\section{Population and prediction tasks}
The existing 39/13/13 development/internal/final partition is retained, with Rudy
outside model fitting. Features are extracted without reference values;
published centers enter only supervised targets and evaluation. Missing inputs
and pipeline failures retain their cohort rows. Global parameter regression is
the primary task. Catastrophic conventional-estimator failure and nominal-alias
reliability are binary quality tasks when class support permits. The selected
positive-only population cannot calibrate astrophysical signal absence.
\section{Models and grouped validation}
For quantity $q$, the target is $y_i=\log(q_i^{\rm ref}/q_i^{\rm classical})$.
The fitted correction is $\widehat q_i=q_i^{\rm classical}\exp(\beta_0+x_i^T\beta)$.
Ridge parameters minimize
\[\sum_i w_i(y_i-\beta_0-x_i^T\beta)^2+\lambda\|\beta\|^2,\]
where each parent star has total weight one. Median imputation and standardization
are fitted inside training folds. Three deterministic outer star folds and three
inner star folds isolate penalty selection from out-of-sample prediction.
The registered grid is $\lambda\in\{1,10,100\}$. Logistic quality models use the
same grouped design; an inner-OOF logit layer calibrates probabilities without
outer test outcomes. Unsupported class/fold fits are withheld.
\section{Uncertainty and practical utility}
The primary loss is absolute fractional error. The paired increment is
$d_i=|\widehat q_i^{B}/q_i^{\rm ref}-1|-|\widehat q_i^{H}/q_i^{\rm ref}-1|$.
Prespecified practical margins are 0.01 for $\nu_{\max}$ and 0.005 for
$\Delta\nu$. Superiority and two-one-sided equivalence tests receive separate
three-test superiority and separate two-regression equivalence corrections. Paired star bootstrap intervals are conditional
on the fitted OOF models; full fitting-selection uncertainty is not represented.
Empirical intervals use inner OOF residual ranks. They are not exact conformal
intervals, and insufficient rank support produces a missing interval. Coverage
uses published reference centers rather than known stellar truth. Quality-model
Brier scores, fixed-bin calibration and threshold decision curves are reported
only with both outcome classes. TESS risk probabilities are withheld under
mission shift. A quality flag is not an oscillation detection.
\section{Current grouped evidence}
\begin{tabular}{lrrrr}
Quantity & Stars & Baseline MAE & Hybrid MAE & Improvement \\
\hline
numax & 9 & 0.02456 & 0.02445 & 0.00011 \\
dnu & 9 & 0.00768 & 0.00807 & -0.00039 \\
\end{tabular}

\par\medskip
These development OOF measurements are a sparse pilot. Unusable stars are not
imputed as successful observations. The source archive includes inherited
pipeline failures; their reasons and compatibility repairs are preserved. The
registered minimum of 20 development stars is a claim-support guard, not a
formal assertion of statistical power. A signed model release must precede the
original Stage 11 holdout opening, after which the first archive is applied once.
\IfFileExists{../figures/stage12/stage12_oof_increment.pdf}{%
\begin{figure}[ht]\centering
\includegraphics[width=\linewidth]{../figures/stage12/stage12_oof_increment.pdf}
\caption{Paired outer-fold errors and conditional bootstrap increments.
The sparse development pilot does not support population complementarity.}\end{figure}}{}
\section{Physical degradation and contribution trajectories}
Kepler Q5 samples are integrated as a piecewise-linear flux curve over target
exposures. Unsupported exposures and long source gaps are rejected. Cadence,
available duration, noise, real sector masks, extra duty loss and contamination
are changed along registered trajectories. Seeds depend on parent star, sector,
source offset and draw; paired methods and degradation levels share random
numbers. Every generated time/flux array and analysis checkpoint is archived.
Contamination dilutes variability before instrumental noise is added. Inherited
Kepler noise is retained. The unit color-response assumption is a sensitivity;
it does not constitute a measured Kepler/TESS passband conversion. The noise
model and real-window coverage remain explicit nuisance limitations.
\section{Actual Rudy sectors and combinations}
The frozen model digest is checked before and during sector processing. All 11
registered SPOC sectors are retained, including quality rejection and
nondetection cases. Actual sectors were previously inspected in Stage 10, so
this is a reproducible case study rather than new blinded test data. Registered
adjacent combinations average PSDs incoherently with exposure weights; they do
not coherently stitch stellar phases across years. Candidate estimates are
retained separately from accepted measurements. Repeated sectors, offsets and
draws still describe one parent star. Conditional degradation quantiles are not
population confidence intervals.
\IfFileExists{../figures/stage12/stage12_transfer_trajectories.pdf}{%
\begin{figure}[ht]\centering
\includegraphics[width=\linewidth]{../figures/stage12/stage12_transfer_trajectories.pdf}
\caption{Rudy cumulative degradation trajectories. Rejected candidates remain
visible; recoverability additionally requires quality acceptance.}\end{figure}}{}
\section{Recalibration, scope and counterexamples}
A separately authorized layer may estimate a star-weighted median log offset
using development simulations from non-Rudy stars. Holdout stars, Rudy and actual
test sectors are forbidden from calibration fitting. Features remain fixed, and
the calibration layer requires evaluation on disjoint simulation stars. No layer
is applied in this snapshot. Missing envelopes, competing aliases, source gaps,
out-of-range features, unavailable classes and mission shift are explicit failure
domains. The current scope decision stops at Rudy transfer limits; no powered
broader TESS sample is established.
\section{Reproducibility and claims}
The release supplies model JSON, feature/task registries, grouped fold audits,
paired results, source-backed figures, degraded light curves, quality reasons,
and a model card. Base inference requires NumPy/SciPy only; mission adapters use
the existing optional astronomy environment. G7's delay-topology drop and G8's
geometry-only interpretation are preserved. Population complementarity requires
locked out-of-sample improvement beyond a practical margin or supported failure
calibration; positive transfer requires compatible actual observations within
uncertainty. Until those conditions are supported, the honest output is a
validation-pending pilot and stated transfer limits.
\bibliographystyle{plain}\bibliography{references}
\end{document}
